\documentclass[11pt]{article}
\usepackage[margin=2.5cm]{geometry}
\usepackage[colorlinks=true, urlcolor=blue, linkcolor=blue]{hyperref}
\usepackage{enumitem}
\usepackage{amsmath}
\usepackage{booktabs}

\title{End-to-End Data \& Experiment Pipeline\\for the SIC Project (multi-hop Semantic Induction Circuits)}
\author{Working plan --- based on Ren et al.\ 2024 and the project proposal}
\date{September 2026}

\begin{document}
\maketitle

\noindent This document lists everything we need --- data, models, tooling, and processing steps --- in the order we should build it, adapted from the experimental setup of
\href{https://arxiv.org/abs/2402.13055}{Ren et al.\ 2024 (Identifying Semantic Induction Heads)} to our research question (distributed Semantic Induction Circuits for multi-hop relations).

\section{Pipeline overview (build in this order)}
\begin{enumerate}[itemsep=2pt]
  \item Choose model suite + checkpoints; verify TransformerLens support.
  \item Generate synthetic multi-hop ICL datasets (gold relations known by construction).
  \item Tokenize and anchor: map every entity to its token indices; fix a multi-token convention.
  \item Run with activation caching; compute per-head Relation Index (RI), reproducing Ren et al.\ on 1-hop as a sanity check.
  \item Generalize RI to circuits; run activation / path patching; measure faithfulness of the sparse subgraph.
  \item Sweep across training checkpoints; separate format-compliance failures from pattern-discovery failures.
\end{enumerate}

\section{Models and checkpoints}
Ren et al.\ used InternLM2-1.8B (24 layers $\times$ 16 heads) for the head analysis, and trained their own 1.1B model on SlimPajama with checkpoints every 200 steps for the ICL-emergence experiments. We do \emph{not} train: we use public checkpoint suites.
\begin{itemize}[itemsep=2pt]
  \item \textbf{Pythia} (EleutherAI): 154 checkpoints per model size, fully public training data order. Repo: \url{https://github.com/EleutherAI/pythia}; models on HF, e.g.\ \href{https://huggingface.co/EleutherAI/pythia-160m}{pythia-160m} \dots\ \href{https://huggingface.co/EleutherAI/pythia-1b}{pythia-1b} (checkpoints via the \texttt{revision} argument, e.g.\ \texttt{revision="step3000"}).
  \item \textbf{OLMo} (AI2): \url{https://github.com/allenai/OLMo}; checkpoints on HF.
  \item \textbf{LMEnt} suite (named in our proposal): locate via \href{https://arxiv.org/list/cs.CL/recent}{arXiv} / \href{https://huggingface.co/models?search=lment}{HF search} and confirm checkpoint availability before committing to it.
  \item \textbf{Tooling}: \href{https://github.com/TransformerLensOrg/TransformerLens}{TransformerLens} --- check the \href{https://transformerlensorg.github.io/TransformerLens/}{supported-models list}; Pythia and OLMo are supported. All caching and patching goes through \texttt{run\_with\_cache} and hook points.
\end{itemize}
\textbf{Decision to make early:} model size vs.\ number of checkpoints vs.\ storage. Attention patterns are $L \times H \times n^2$ per prompt \emph{per checkpoint}; compute RI on the fly and store only per-head statistics, never raw patterns for the full sweep.

\section{Datasets}
\subsection{What Ren et al.\ used (our 1-hop reference)}
\begin{itemize}[itemsep=2pt]
  \item \href{https://github.com/rikdz/GraphWriter}{AGENDA} (Koncel-Kedziorski et al.\ 2019): 1{,}000 test samples; each = a scientific abstract + a knowledge graph of entities and typed relations, \emph{anchored in the text}. Used for knowledge-graph relations.
  \item Syntactic dependencies: sentences split from the same paragraphs, parsed with \href{https://spacy.io}{spaCy}. Note: these gold labels are themselves model output (parser noise) --- a caveat they inherit and we avoid.
  \item ICL-levels tasks: few-shot classification prompts (their Table~1); format compliance = the forced single generated token has the right \emph{type} (e.g.\ a digit), pattern discovery = the label is \emph{correct}; class-balanced shots.
\end{itemize}

\subsection{What we build (multi-hop, synthetic)}
Generated from templates so that gold relations are known by construction --- no parser, no annotation noise.
\begin{itemize}[itemsep=2pt]
  \item \textbf{Template sources}: \href{https://github.com/facebookresearch/clutrr}{CLUTRR} (compositional kinship reasoning; built for exactly this) and \href{https://research.facebook.com/downloads/babi/}{bAbI} tasks. Per the proposal, replace all real nouns with fictitious entities so the model cannot answer from parametric memory.
  \item \textbf{Each sample must carry}: the prompt text; the list of gold triples $(\text{head},\ \text{relation},\ \text{tail})$; for multi-hop, the chain (e.g.\ $A\!\subset\!B$, $B\!\subset\!C \Rightarrow A\!\subset\!C$) with the \emph{bridge entity} marked; character-level spans of every entity occurrence.
  \item \textbf{Matched difficulty tiers}: 1-hop (reproduces Ren et al.'s setting), 2-hop, optionally 3-hop --- same entities, same format, so hop count is the only manipulated variable.
  \item \textbf{Controls to generate alongside} (needed later for patching and for RI baselines):
    corrupted prompts (bridge entity replaced --- breaks the chain while preserving format);
    shuffled-relation prompts; and Ren et al.'s trick of a \emph{fake tail} baseline (they re-computed RI with the tail forced to an arbitrary token --- e.g.\ the 10th token --- as a null distribution).
  \item \textbf{Anti-confound rule}: the query pair must never appear verbatim as an in-context example --- otherwise a plain (non-semantic) induction head solves the task by literal copying and inflates our results.
  \item \textbf{Class balance}: when shots $\geq$ number of answer classes, include at least one example per class (Ren et al.'s protocol), so that format accuracy is well-defined.
\end{itemize}

\section{Tokenization and anchoring}
The single most error-prone step. Annotations live on words/entities; the model lives on sub-word tokens.
\begin{enumerate}[itemsep=2pt]
  \item Tokenize each prompt with the \emph{suite's own tokenizer}; build a map from every entity occurrence (character span) to a contiguous range of token indices.
  \item Fix a convention for multi-token entities (fictitious names will split into several tokens). Ren et al.\ operate at the token level; for spans, the common convention is to use the \emph{last} token of the entity as its representative (where its identity is complete). Document the choice and test robustness to it.
  \item Store, per sample: \texttt{(prompt, token\_ids, \{entity $\to$ [token idx range]\}, gold triples, chain structure)} --- e.g.\ one JSONL line per sample.
\end{enumerate}

\section{Metrics}
\subsection{Relation Index (Ren et al., \S3) --- reproduce first on 1-hop}
Two stages, per head $h$, per gold triple with head token $t_s$, tail token $t_o$, evaluated at current tokens $t_j$ with $j \geq \max(s,o)$:
\begin{enumerate}[itemsep=2pt]
  \item \textbf{QK stage (attention selection)}: the head must attend from $t_j$ to $t_s$ dominantly: $s = \arg\max_k A^h_{j,k}$ and $A^h_{j,s} / \max_{k \neq s} A^h_{j,k} > \tau$ (they use $\tau = 2.2$; see their Appendix~A). Otherwise the head is skipped for this triple.
  \item \textbf{OV stage (tail promotion)}: project the head's OV output to the vocabulary, $p^{h,j} = \mathrm{softmax}(x_j W^h_{OV} W_U)$; keep only above-mean probabilities $q^{h,j}_{t_k} = \max(0,\ p^{h,j}_{t_k} - \mathbb{E}_{k'}[p^{h,j}_{t_{k'}}])$; the score is the tail's share $a^{h,j}_T = q^{h,j}_{t_o} / \sum_k q^{h,j}_{t_k}$. Average over $j$ and over triples.
\end{enumerate}
\textbf{Our extension (core novelty)}: generalize RI from a single head to a \emph{circuit}: for 2-hop, no single head needs to link $A$ to $C$; instead head $h_1$ may resolve $A\!\to\!B$ and write it to the residual stream, and $h_2$ read it and resolve $\to C$ (composition). The circuit-level RI must therefore score (i) per-edge behavior and (ii) the causal link between the heads --- which is what path patching measures (below), so the metric and the intervention should be designed together.

\subsection{ICL levels (for the checkpoint sweep)}
Follow Ren et al.\ \S4: \textbf{loss reduction} --- interval-averaged version of Olsson et al.'s score (mean loss over tokens $j\ldots i{+}j$ minus mean over $0\ldots i$); \textbf{format compliance} --- forced single-token generation, check the token's type; \textbf{pattern discovery} --- label accuracy. Report all three per checkpoint so that ``no SIC yet'' is never confused with ``model cannot even comply with the format yet''.

\section{Causal experiments}
To learn/implement (in this order):
\begin{enumerate}[itemsep=2pt]
  \item \textbf{Activation patching} (node-level): replace a head's/MLP's activation with its value from the corrupted prompt; measure the drop in the correct-answer probability. Primer: \href{https://arxiv.org/abs/2404.15255}{Heimersheim \& Nanda 2024, ``How to use and interpret activation patching''}; origins: \href{https://arxiv.org/abs/2202.05262}{ROME / causal tracing}.
  \item \textbf{Path patching} (edge-level): patch only along the path from component $X$ to component $Y$ --- this is what substantiates ``$h_1$ feeds $h_2$'', i.e.\ the circuit claim. Canonical use: \href{https://arxiv.org/abs/2211.00593}{Wang et al.\ 2022 (IOI)} --- near-mandatory reading; our faithfulness protocol is modeled on it.
  \item \textbf{Faithfulness} (proposal \S3): run the model with only the identified sparse subgraph active (everything else mean-ablated or frozen); compare correct-answer probability to the full model. Define the recovery threshold \emph{in advance}. Rigor ceiling for this kind of claim: \href{https://www.alignmentforum.org/posts/JvZhhzycHu2Yd57RN/causal-scrubbing-a-method-for-rigorously-testing}{causal scrubbing} (skim).
\end{enumerate}

\section{Sanity checks and baselines}
\begin{itemize}[itemsep=2pt]
  \item Reproduce Ren et al.\ qualitatively on 1-hop relations before any 2-hop work (small scale: a few relations, a few hundred samples).
  \item Fake-tail / random-tail RI baseline (their trick) as the null distribution for every RI heatmap.
  \item Plain induction-head baseline: measure prefix-matching and copying scores (Olsson et al.) on the same heads, to show SIC heads are not just literal-copy heads --- and vice versa.
  \item Random / earliest-checkpoint model as a floor for every metric.
  \item Verify the anti-confound rule held (no verbatim query repetition) by scanning generated prompts.
\end{itemize}

\section{Storage and compute notes}
$\sim$100--1{,}000 prompts $\times$ dozens of checkpoints: cache per-head \emph{scores}, not raw attention ($L\,H\,n^2$ floats per prompt adds up fast); batch prompts per checkpoint load (loading a checkpoint dominates wall-time); models $\leq$1B fit the studentkillable Slurm partition per the proposal.

\section{Key resources (one place)}
\begin{itemize}[itemsep=1pt]
  \item Ren et al.\ 2024 (SIH): \url{https://arxiv.org/abs/2402.13055}
  \item Olsson et al.\ 2022 (induction heads): \url{https://transformer-circuits.pub/2022/in-context-learning-and-induction-heads/index.html}
  \item Elhage et al.\ 2021 (framework): \url{https://transformer-circuits.pub/2021/framework/index.html}
  \item Wang et al.\ 2022 (IOI, path patching): \url{https://arxiv.org/abs/2211.00593}
  \item Activation patching guide: \url{https://arxiv.org/abs/2404.15255}
  \item TransformerLens: \url{https://github.com/TransformerLensOrg/TransformerLens}
  \item Pythia: \url{https://github.com/EleutherAI/pythia} \quad OLMo: \url{https://github.com/allenai/OLMo}
  \item CLUTRR: \url{https://github.com/facebookresearch/clutrr} \quad bAbI: \url{https://research.facebook.com/downloads/babi/}
\end{itemize}

\end{document}
